\documentclass[12pt]{article}
\usepackage[utf8]{inputenc}
\usepackage[margin=1in]{geometry}
\usepackage{tabularx}
\usepackage{booktabs}
\usepackage{hyperref}

\title{Assignment 2: Module Interface \& Header Design}
\author{Autonomous Fire Detection and Suppression Rover}
\date{March 28, 2026}

\begin{document}

% --- 3.1 COVER PAGE ---
\begin{titlepage}
    \centering
    \vspace*{0.5in}
    \Huge \textbf{CSE 396 - Computer Engineering Project} \\
    \vspace{0.4in}
    \huge \textbf{Assignment 2: Module Interface \& Header Design} \\
    \vspace{0.6in}
    \Large \textbf{Project Title:} \\
    Autonomous Fire Detection and Suppression Rover \\
    \vspace{0.6in}
    \Large \textbf{Team Members \& Primary Assignments:} \\
    \normalsize
    \vspace{0.2in}
    \begin{tabularx}{\textwidth}{l l X}
        \textbf{Name} & \textbf{Student ID} & \textbf{Assigned Submodules} \\ \midrule
        Görkem UYSAL & 230104004174 & MOD-07, MOD-08 \\
        Cemal GÜMÜŞ & 210104004304 & MOD-01, MOD-03 \\
        Melih KAZANCI & 230104004909 & MOD-01, MOD-03 \\
        Ahmet Halil YAMLI & [ID] & MOD-04, MOD-05 \\
        Derya UYSAL & 220104004045 & MOD-04, MOD-05 \\
        Mehmet Alp ATAY & 220104004020 & MOD-02, MOD-06 \\
        Adil Mert ERGÖRÜN & 220104004048 & MOD-02, MOD-06 \\
        Ayşe Feyza SERBEST & 220104004052 & MOD-07, MOD-08 \\
    \end{tabularx}
    \vfill
    \Large \textbf{Date:} March 28, 2026
\end{titlepage}

% --- 3.2 MODULE-STUDENT MATCHING TABLE ---
\section{Module-Student Matching Table}
The following table provides the formal mapping of team members to the technical submodules defined for the rover system.

\begin{table}[h]
\centering
\small
\begin{tabularx}{\textwidth}{l l l X}
\toprule
\textbf{ID} & \textbf{Module Name} & \textbf{Responsible Student(s)} & \textbf{Brief Role} \\ \midrule
MOD-01 & Flame Sensing HW & Melih K., Cemal G. & Sensor mounting, wiring, and noise reduction. \\
MOD-02 & Flame Detect. SW & Mehmet Alp A., Adil Mert E. & Signal filtering and direction estimation. \\
MOD-03 & Locomotion HW & Melih K., Cemal G. & Chassis assembly and motor driver wiring. \\
MOD-04 & Motion Control SW & Derya U., Ahmet Halil Y. & Differential drive and speed control logic. \\
MOD-05 & Dist. Sensing & Derya U., Ahmet Halil Y. & HC-SR04 integration and thresholding. \\
MOD-06 & Suppression HW & Mehmet Alp A., Adil Mert E. & Water pump, relay, and tank integration. \\
MOD-07 & Autonomy SW & Görkem U., Ayşe Feyza S. & FSM design and decision logic coordination. \\
MOD-08 & Integration/Test & Görkem U., Ayşe Feyza S. & End-to-end verification and debugging. \\ \bottomrule
\end{tabularx}
\end{table}

% --- 3.3 SYSTEM / MODULE OVERVIEW ---
\section{System / Module Overview}
The rover utilizes a distributed architecture where hardware-facing modules (MOD-01, 03, 05, 06) interact with software processing modules (MOD-02, 04, 07) to achieve autonomous suppression.

\subsection{Module Purpose Summaries}
\begin{itemize}
    \item \textbf{MOD-01 \& MOD-02:} Together, these modules handle the "Detect" phase of the FSM by providing a filtered direction vector toward the flame.
    \item \textbf{MOD-03 \& MOD-04:} These handle the physical movement of the robot, translating high-level commands (e.g., "Turn Left") into PWM signals for the motors.
    \item \textbf{MOD-05:} This module provides critical safety feedback, ensuring the rover stops within the required distance to prevent collision with the fire source.
    \item \textbf{MOD-06:} Manages the final suppression stage, activating the relay for the water pump once the target is reached.
    \item \textbf{MOD-07 \& MOD-08:} These oversee the system states and ensure all submodules interact correctly during full-cycle operation.
\end{itemize}

% --- 3.5 INTER-MODULE COMMUNICATION DESCRIPTION ---
\section{Inter-Module Communication Description}
The communication architecture relies on structured data exchange between hardware drivers and control logic.

\begin{table}[h]
\centering
\footnotesize
\begin{tabularx}{\textwidth}{l l l l l X}
\toprule
\textbf{From} & \textbf{To} & \textbf{Data / Signal} & \textbf{Type} & \textbf{Dir.} & \textbf{Notes} \\ \midrule
MOD-01 & MOD-02 & \texttt{raw\_adc\_values} & \texttt{uint16\_t[]} & $01 \to 02$ & Raw IR intensity per sensor. \\
MOD-02 & MOD-07 & \texttt{flame\_data\_t} & \texttt{struct} & $02 \to 07$ & Direction enum and intensity proxy for distance. \\
MOD-05 & MOD-07 & \texttt{distance\_mm} & \texttt{uint16\_t} & $05 \to 07$ & Target distance; 20Hz update. \\
MOD-07 & MOD-04 & \texttt{motion\_cmd\_t} & \texttt{struct} & $07 \to 04$ & Includes direction and speed. \\
MOD-07 & MOD-06 & \texttt{pump\_active} & \texttt{bool} & $07 \to 06$ & Trigger for suppression relay. \\ \bottomrule
\end{tabularx}
\end{table}

% --- 3.6 KNOWN RISKS ---
\section{Known Risks \& Open Questions}
\begin{enumerate}
    \item \textbf{Communication Latency:} Frequent polling of the ultrasonic sensor (MOD-05) may block the FSM (MOD-07) execution if not implemented using non-blocking interrupts.
    \item \textbf{Module Overlap:} Since several members span multiple modules (e.g., MOD-04 and MOD-05), maintaining strict interface boundaries in code will be critical to prevent tight coupling.
\end{enumerate}

\end{document}